\chapter{Supplementary Material}
\label{app:supplementary}

\section{Repository Structure}
\label{app:repo-structure}

The source code is organised as a monorepo with the following top-level layout.

\begin{lstlisting}
DACNTT_VideoStreaming/
  frontend/            React single page application (Vite, HLS.js)
  backend/             Node.js Express API and test suite
  transcoder/          FFmpeg worker and Docker image
  infrastructure/
    modules/           Eleven reusable Terraform modules
    environments/      dev and prod instantiations
  .github/workflows/   Seven CI/CD workflows
  scripts/             Benchmarking and load testing scripts
  docs/                Architecture diagrams, cost analysis, measured results
  report/              LaTeX sources of this report
\end{lstlisting}

\section{FFmpeg Transcoding Command}
\label{app:ffmpeg}

The following invocation produces all three renditions from one read of the source. A single
input is followed by three output blocks, each carrying its own scaling filter and encoder
settings, so FFmpeg decodes the source once and feeds all three filter chains from it. The
second and third blocks are abbreviated here; in the source they repeat the first block in
full. The master playlist is not written by FFmpeg: the worker generates it afterwards, which is
why no \texttt{-master\_pl\_name} argument appears. It does so from the encoded output rather
than from the settings above, measuring each rendition's peak segment bit rate from the segment
sizes and their \texttt{\#EXTINF} durations, and reading the codec string for each variant out of
the sequence parameter set of its first segment.

The keyframe interval is derived from the source rather than assumed. The worker probes the
first video stream and reads \texttt{avg\_frame\_rate}, preferring it over
\texttt{r\_frame\_rate} because the latter is a base rate FFmpeg infers from the
timing grid and can report values such as 1000/1 for variable-frame-rate recordings. The group
of pictures is then \texttt{floor(fps} $\times$ \texttt{6)} frames, floored rather than rounded so that
it never exceeds the segment length: at 29.97~fps, rounding would give 180 frames, or 6.006
seconds, and the segment boundary would drift past the point where a keyframe exists. The value
shown below, 144, is what a 24~fps source produces.

Two further arguments keep the boundaries aligned. Scene detection is disabled, because letting
the encoder place keyframes at scene changes would allow the three renditions to cut at
different points, and RFC~8216 requires that matching content in variant streams carry matching
timestamps. \texttt{-force\_key\_frames} then pins a keyframe to every sixth second of
media time; unlike \texttt{-g}, which only bounds the interval from above, it is absolute and
independent of the detected frame rate, so a misdetected source still yields aligned segments.

The H.264 profile is pinned but the level is not. Each rendition carries a different frame size,
and the lowest conforming level therefore differs between them, so the level is left for the
encoder to derive from the resolution rather than fixed in the command.

\begin{lstlisting}[language=bash]
ffmpeg -i input.mp4 -hide_banner -loglevel warning -y \
  -vf "scale=640:360:force_original_aspect_ratio=decrease, \
         pad=640:360:(ow-iw)/2:(oh-ih)/2" \
  -c:v libx264 -b:v 400k  -maxrate 500k  -bufsize 800k  \
  -preset fast -profile:v main \
  -g 144 -keyint_min 144 -sc_threshold 0 \
  -force_key_frames "expr:gte(t,n_forced*6)" \
  -c:a aac -b:a 64k  -ar 44100 -ac 2 \
  -f hls -hls_time 6 -hls_list_size 0 \
  -hls_segment_filename 360p/segment_%03d.ts  360p/playlist.m3u8 \
  \
  -vf "scale=1280:720:..."  -c:v libx264 -b:v 1500k -maxrate 2000k -bufsize 3000k \
  -c:a aac -b:a 128k -f hls -hls_time 6 \
  -hls_segment_filename 720p/segment_%03d.ts  720p/playlist.m3u8 \
  \
  -vf "scale=1920:1080:..." -c:v libx264 -b:v 4000k -maxrate 5000k -bufsize 8000k \
  -c:a aac -b:a 192k -f hls -hls_time 6 \
  -hls_segment_filename 1080p/segment_%03d.ts 1080p/playlist.m3u8
\end{lstlisting}

\section{Signed Cookie Custom Policy}
\label{app:signed-policy}

The policy signed by the backend before issuing playback cookies. The wildcard in the resource
path covers the manifest and every segment of one specific video, and nothing beyond it.

\begin{lstlisting}[language=JavaScript]
{
  "Statement": [
    {
      "Resource": "https://cdn.zelostech.site/videos/<videoId>/*",
      "Condition": {
        "DateLessThan": { "AWS:EpochTime": 1786300347 }
      }
    }
  ]
}
\end{lstlisting}

\section{Reproducing the Measurements}
\label{app:reproduce}

The experimental results reported in Chapter~\ref{chap:evaluation} are produced by the
scripts listed below. Detailed instructions, including the environment variables required and
the order in which the measurements must be taken, are given in
\texttt{docs/results/README.md} within the repository.

Every table of measurements in Chapter~\ref{chap:evaluation} is listed below against the script
that produced it and the file its figures were read from, so that any individual number can be
traced back to its source. Paths are relative to the repository root, and the result files are
those under \path{docs/results/}.

\begin{table}[H]
\centering
\caption{Provenance of each measurement table}
\label{tab:provenance}
\footnotesize
\begin{tabularx}{\textwidth}{l p{0.30\textwidth} X}
\toprule
\textbf{TABLE} & \textbf{SCRIPT} & \textbf{RESULT FILE} \\
\midrule
\ref{tab:ttff-results}       & \path{scripts/benchmark-qoe.js}  & \path{qoe-ttff.json} \\
\ref{tab:ttff-cloudfront}    & \path{scripts/benchmark-qoe.js}  & \path{qoe-ttff-cloudfront.json} \\
\ref{tab:ttff-multiregion}   & \path{scripts/benchmark-qoe.js}, deployed as a Lambda to each region & \path{qoe-ttff-multiregion.json} \\
\ref{tab:qoe-playback}       & \path{scripts/qoe/collect.js}    & \path{qoe-playback-<profile>.json} \\
\ref{tab:qoe-bandwidth}      & Segment sizes and \texttt{\#EXTINF} durations read from the processed bucket & --- \\
\ref{tab:qoe-before-after}   & \path{scripts/qoe/collect.js}, run before and after the correction & \path{qoe-playback-<profile>.json} \\
\ref{tab:transcode-results}  & \path{scripts/measure-transcode.js} & \path{transcode-timing.json} \\
\ref{tab:load-results}       & \path{scripts/k6-load-test.js}   & \path{k6-summary.json}, \path{stress-test-concurrency.json} \\
\ref{tab:load-crossvalidation} & \path{scripts/node-load-test.js} & \path{node-load-test.json} \\
\ref{tab:drain-rate}         & \path{scripts/k6-load-test.js}   & \path{drain-rate-real-payload.json}, \path{stress-test-concurrency.json} \\
\ref{tab:cost-comparison}    & Published AWS list prices, computed analytically & --- \\
\bottomrule
\end{tabularx}
\end{table}

\section{Deploying the Infrastructure}
\label{app:deploy}

The complete AWS environment is provisioned from a single directory.

\begin{lstlisting}[language=bash]
cd infrastructure/environments/dev
terraform init
terraform apply
\end{lstlisting}

Enabling signed cookies additionally requires generating an RSA key pair, supplying the public
key to Terraform, and configuring the backend with the corresponding private key and the key
pair identifier reported by \texttt{terraform output}. The full procedure is documented in
\texttt{docs/PRIVATE\_VIDEO\_SIGNED\_COOKIES.md}.
